% !PRINTED 158
\weberpage{158}

%<<CONT
schiedenen Formen gibt, die nur mühsam aufeinander zurückführbar
 sind. Zunächst erhalten wir aus der Definition (1),
wenn wir den Parameter jetzt mit $a$ bezeichnen:

\weq{6}{\begin{aligned}
&\Pi(u+v,a)-\Pi(u,a)-\Pi(v,a)\\
&\qquad = \dfrac{1}{2}\log\dfrac{\Theta(u+v-a)}{\Theta(u+v+a)}\,\dfrac{\Theta(u+a)\,\Theta(v+a)}{\Theta(u-a)\,\Theta(v-a)},
\end{aligned}}

und es handelt sich noch darum, den unter dem Logarithmus
stehenden Ausdruck durch elliptische Funktionen darzustellen.

Dies geschieht zunächst leicht, wie an der erwähnten Stelle
der Fundamenta, mittels der Formel [§ 43, (11)]:

\[
\Theta^2(0)\,\Theta(u+v)\,\Theta(u-v) = \Theta^2(u)\,\Theta^2(v)\,[1-\varkappa^2\,\mathrm{sn}^2u\,\mathrm{sn}^2v],
\]

wenn man darin für $u,v$ setzt $u\pm a$, $v\pm a$, dann $\pm a$,
$u+v\pm a$.

Man erhält so:

\[
\begin{aligned}
\Theta^2(0)\,\Theta(u+v\pm 2a)\,\Theta(u-v) &= \Theta^2(u\pm a)\,\Theta^2(v\pm a)\,[1-\varkappa^2\,\mathrm{sn}^2(u\pm a)\,\mathrm{sn}^2(v\pm a)],\\
\Theta^2(0)\,\Theta(u+v\pm 2a)\,\Theta(u+v) &= \Theta^2(a)\,\Theta^2(u+v\pm a)\,[1-\varkappa^2\,\mathrm{sn}^2a\,\mathrm{sn}^2(u+v\pm a)],
\end{aligned}
\]

woraus folgt:

\weq{7}{\dfrac{\Theta(u+v-a)\,\Theta(u+a)\,\Theta(v+a)}{\Theta(u+v+a)\,\Theta(u-a)\,\Theta(v-a)} = \sqrt{\dfrac{[1-\varkappa^2\,\mathrm{sn}^2(u-a)\,\mathrm{sn}^2(v-a)]\,[1-\varkappa^2\,\mathrm{sn}^2a\,\mathrm{sn}^2(u+v+a)]}{[1-\varkappa^2\,\mathrm{sn}^2(u+a)\,\mathrm{sn}^2(v+a)]\,[1-\varkappa^2\,\mathrm{sn}^2a\,\mathrm{sn}^2(u+v-a)]}}.}

Einen zweiten Ausdruck erhält man aus der Additionsformel
(13), § 22, die man so darstellen kann:

\[
\Theta(0)\,\Theta(u\pm a)\,\Theta(v\pm a)\,\Theta(u+v) = \Theta(u)\,\Theta(v)\,\Theta(a)\,\Theta(u+v\pm a)\,[1\pm\varkappa^2\,\mathrm{sn}\,a\,\mathrm{sn}\,u\,\mathrm{sn}\,v\,\mathrm{sn}(u+v\pm a)],
\]

woraus durch Division:

\weq{8}{\dfrac{\Theta(u+v-a)\,\Theta(u+a)\,\Theta(v+a)}{\Theta(u+v+a)\,\Theta(u-a)\,\Theta(v-a)} = \dfrac{1+\varkappa^2\,\mathrm{sn}\,a\,\mathrm{sn}\,u\,\mathrm{sn}\,v\,\mathrm{sn}(u+v+a)}{1-\varkappa^2\,\mathrm{sn}\,a\,\mathrm{sn}\,u\,\mathrm{sn}\,v\,\mathrm{sn}(u+v-a)}.}

\gesperrt{Jacobi} gibt noch einen dritten Ausdruck für die
$\Theta$-Quotienten in (6):

\weq{9}{\dfrac{\left\{1-\varkappa^2\,\mathrm{sn}^2\!\left(\dfrac{u-v}{2}\right)\mathrm{sn}^2\!\left(\dfrac{u+v}{2}+a\right)\right\}\left\{1-\varkappa^2\,\mathrm{sn}^2\!\left(\dfrac{u+v}{2}\right)\mathrm{sn}^2\!\left(\dfrac{u+v}{2}-a\right)\right\}}{\left\{1-\varkappa^2\,\mathrm{sn}^2\!\left(\dfrac{u-v}{2}\right)\mathrm{sn}^2\!\left(\dfrac{u+v}{2}-a\right)\right\}\left\{1-\varkappa^2\,\mathrm{sn}^2\!\left(\dfrac{u+v}{2}\right)\mathrm{sn}^2\!\left(\dfrac{u+v}{2}+a\right)\right\}}.}
